\documentclass[a4paper, final]{article}

\usepackage[russian]{babel}
\usepackage{fontspec}
\setmainfont{Times New Roman}
\newfontfamily\cyrillicfonttt{Courier New}
\usepackage[14pt]{extsizes}
\usepackage[left=20mm, top=20mm, right=20mm, bottom=20mm, footskip=10mm]{geometry}
\usepackage{ragged2e}
\usepackage{indentfirst}
\usepackage{graphicx}
\usepackage{caption}
\usepackage{xcolor}
\usepackage{listings}
\usepackage{booktabs}
\usepackage{hyperref}
\usepackage{amsmath}
\usepackage{amssymb}
\usepackage{array}
\usepackage{multirow}

\definecolor{link_color}{HTML}{0000EE}
\hypersetup{colorlinks=true,linkcolor=link_color,urlcolor=link_color,citecolor=link_color}

\lstset{
    language=Python,
    basicstyle=\cyrillicfonttt\small,
    numbers=left,
    numberstyle=\tiny\color{gray},
    keywordstyle=\color{blue}\bfseries,
    commentstyle=\color{gray}\itshape,
    stringstyle=\color{red},
    frame=single,
    captionpos=b,
    breaklines=true,
    breakatwhitespace=false,
    showstringspaces=false,
    keepspaces=true,
    columns=flexible,
    fontadjust=true,
    basewidth={0.5em,0.45em},
    tabsize=4,
}

\begin{document}

{
\thispagestyle{empty}
\begin{center}
    \hfill\break\hfill\break
    МИНИСТЕРСТВО НАУКИ И ВЫСШЕГО ОБРАЗОВАНИЯ РОССИЙСКОЙ ФЕДЕРАЦИИ\\[5pt]
    Федеральное государственное автономное образовательное учреждение высшего образования
    \guillemotleft Санкт-Петербургский политехнический университет Петра Великого\guillemotright\\[10pt]
    Институт компьютерных наук и кибербезопасности\\[5pt]
    Направление: 02.03.01 Математика и компьютерные науки\\[30pt]
    {\large Отчет по дисциплине: \guillemotleft Вычислительная математика\guillemotright}\\[30pt]
    \LARGE\bf \guillemotleft Лабораторная работа 5\guillemotright.
\end{center}
\vspace{110pt}
\begin{tabular*}{460pt}{@{\extracolsep{\fill}} l r}
    Студент,\tabularnewline группы 5130201/40003 \hspace{50pt} & \line(1,0){100}\quad Четвергов И.\,С.\\
    & \\
    Руководитель,\tabularnewline Преподаватель & \line(1,0){121.5}\quad Пак В.\,Г.\\
    & \\ & \\ & \\
    & \guillemotleft\line(1,0){30}\guillemotright\line(1,0){100}\,2026\,г.
\end{tabular*}\\
\vfill
\begin{center}Санкт-Петербург, 2026\end{center}
\newpage
}

\tableofcontents
\newpage

\section{Постановка задачи}

\textbf{Тема:} задача о трёхмерной транспортировке (транспортная задача).

\textbf{Цель:} изучить методы решения транспортных задач: метод минимального элемента,
метод Фогеля, решение через симплекс-метод.

\textbf{Задание:}
\begin{enumerate}
    \item Построить начальное допустимое решение методом минимального элемента.
    \item Применить метод Фогеля для получения лучшего начального решения.
    \item Оценить стоимость перевозок для каждого решения.
    \item Провести анализ результатов.
\end{enumerate}

\subsection{Исходные данные}

\textbf{Вариант 56}. Матрица интервалов (стоимость транспортировки $c_{ij}$ за единицу груза):

\[
C = \begin{pmatrix}
0{,}07 & 0{,}09 & 0{,}12 & 0{,}10 \\
0{,}11 & 0{,}06 & 0{,}08 & 0{,}13 \\
0{,}10 & 0{,}14 & 0{,}07 & 0{,}09 \\
0{,}12 & 0{,}08 & 0{,}11 & 0{,}08
\end{pmatrix}.
\]

Матрица предложения поставщиков:
\[
a = \begin{pmatrix} 140 \\ 125 \\ 160 \\ 130 \end{pmatrix}, \quad \text{всего:} \quad \sum a_i = 555.
\]

Матрица спроса потребителей:
\[
b = \begin{pmatrix} 130 \\ 145 \\ 155 \\ 125 \end{pmatrix}, \quad \text{всего:} \quad \sum b_j = 555.
\]

Систем�� сбалансирована: $\sum a_i = \sum b_j = 555$. Задача имеет оптимальное решение.

\newpage

\section{Задание 1. Метод минимального элемента}

\subsection{Теоретические сведения}

Метод минимального элемента~--- жадный алгоритм построения начального допустимого
базисного решения транспортной задачи.

\textbf{Алгоритм:}
\begin{enumerate}
    \item Найти клетку $(i,j)$ с минимальной стоимостью в матрице $C$.
    \item Выполнить перевозку объёма $x_{ij} = \min(a_i, b_j)$.
    \item Обнулить соответствующую строку или столбец (что исчерпано раньше).
    \item Повторить до полного распределения груза.
\end{enumerate}

\subsection{Пошаговое решение}

\textbf{Шаг 1.} Минимум $C$: элемент $c_{23} = 0{,}06$ (строка 2, столбец 3).
Перевозим $x_{23} = \min(125, 155) = 125$. Остаток спроса: $b_3 = 155 - 125 = 30$.
Поставщик 2 исчерпан, исключаем строку 2.

\textbf{Шаг 2.} Минимум в оставшейся таблице: $c_{34} = 0{,}09$ (строка 3, столбец 4).
Перевозим $x_{34} = \min(160, 125) = 125$. Остаток предложения: $a_3 = 160 - 125 = 35$.
Потребитель 4 удовлетворен, исключаем столбец 4.

\textbf{Шаг 3.} Минимум: $c_{31} = 0{,}10$ (строка 3, столбец 1).
Перевозим $x_{31} = \min(35, 130) = 35$. Остаток спроса: $b_1 = 130 - 35 = 95$.
Поставщик 3 исчерпан, исключаем строку 3.

\textbf{Шаг 4.} Минимум: $c_{13} = 0{,}12$ (строка 1, столбец 3).
Перевозим $x_{13} = \min(140, 30) = 30$. Остаток предложения: $a_1 = 140 - 30 = 110$.
Потребитель 3 удовлетворен, исключаем столбец 3.

\textbf{Шаг 5.} Минимум: $c_{11} = 0{,}07$ (строка 1, столбец 1).
Перевозим $x_{11} = \min(110, 95) = 95$. Остаток спроса: $b_1 = 95 - 95 = 0$.
Потребитель 1 удовлетворен, исключаем столбец 1.

\textbf{Шаг 6.} Остаток: $a_1 = 15$, $b_2 = 145$. Перевозим $x_{12} = 15$.
Остаток спроса: $b_2 = 145 - 15 = 130$, но поставщик 1 исчерпан.

\textbf{Шаг 7.} Остаток: $a_4 = 130$, $b_2 = 130$. Перевозим $x_{42} = 130$.
Все спрос и предложение распределены.

\subsection{Матрица перевозок (метод минимального элемента)}

\begin{table}[h!]
\centering
\caption{Распределение грузов методом минимального элемента}
\label{tab:min_elem}
\renewcommand{\arraystretch}{1.3}
\begin{tabular}{|c|c|c|c|c|c|}
\hline
 & $j=1$ & $j=2$ & $j=3$ & $j=4$ & $a_i$ \\
\hline
$i=1$ & $95$ & $15$ & $30$ & $0$ & $140$ \\
      & $(0{,}07)$ & $(0{,}09)$ & $(0{,}12)$ &  &  \\
\hline
$i=2$ & $0$ & $0$ & $125$ & $0$ & $125$ \\
      &  &  & $(0{,}06)$ &  &  \\
\hline
$i=3$ & $35$ & $0$ & $0$ & $125$ & $160$ \\
      & $(0{,}10)$ &  &  & $(0{,}09)$ &  \\
\hline
$i=4$ & $0$ & $130$ & $0$ & $0$ & $130$ \\
      &  & $(0{,}08)$ &  &  &  \\
\hline
$b_j$ & $130$ & $145$ & $155$ & $125$ & $555$ \\
\hline
\end{tabular}
\end{table}

\subsection{Расчёт стоимости}

Базисные переменные (ненулевые клетки):
\begin{align*}
x_{11} &= 95, \quad c_{11} = 0{,}07, \quad \text{стоимость} = 95 \cdot 0{,}07 = 6{,}65,\\
x_{12} &= 15, \quad c_{12} = 0{,}09, \quad \text{стоимость} = 15 \cdot 0{,}09 = 1{,}35,\\
x_{13} &= 30, \quad c_{13} = 0{,}12, \quad \text{стоимость} = 30 \cdot 0{,}12 = 3{,}60,\\
x_{23} &= 125, \quad c_{23} = 0{,}06, \quad \text{стоимость} = 125 \cdot 0{,}06 = 7{,}50,\\
x_{31} &= 35, \quad c_{31} = 0{,}10, \quad \text{стоимость} = 35 \cdot 0{,}10 = 3{,}50,\\
x_{34} &= 125, \quad c_{34} = 0{,}09, \quad \text{стоимость} = 125 \cdot 0{,}09 = 11{,}25,\\
x_{42} &= 130, \quad c_{42} = 0{,}08, \quad \text{стоимость} = 130 \cdot 0{,}08 = 10{,}40.
\end{align*}

\textbf{Общая стоимость метода минимального элемента:}
\[
Z_{\text{МЭ}} = 6{,}65 + 1{,}35 + 3{,}60 + 7{,}50 + 3{,}50 + 11{,}25 + 10{,}40 = \mathbf{44{,}25}.
\]

Число базисных переменных: $7 = m + n - 1 = 4 + 4 - 1 = 7$. ✓

\newpage

\section{Задание 2. Метод Фогеля}

\subsection{Теоретические сведения}

Метод Фогеля~--- более сложный алгоритм, учитывающий штраф за неиспользование
дешёвого маршрута. На каждом шаге рассчитывается штраф (разница между минимальной
и второй минимальной стоимостью) по каждой строке и столбцу. Затем выбирается
маршрут в строке или столбце с максимальным штрафом.

\subsection{Пошаговое решение}

\textbf{Итерация 1.}

Штрафы по строкам (разница между минимум и второй минимум):
\begin{align*}
\text{Строка 1:} \quad &\text{мин} = 0{,}07, \quad \text{второй} = 0{,}09, \quad \text{штраф} = 0{,}02,\\
\text{Строка 2:} \quad &\text{мин} = 0{,}06, \quad \text{второй} = 0{,}08, \quad \text{штраф} = 0{,}02,\\
\text{Строка 3:} \quad &\text{мин} = 0{,}07, \quad \text{второй} = 0{,}09, \quad \text{штраф} = 0{,}02,\\
\text{Строка 4:} \quad &\text{мин} = 0{,}08, \quad \text{второй} = 0{,}11, \quad \text{штраф} = 0{,}03.
\end{align*}

Штрафы по столбцам:
\begin{align*}
\text{Столбец 1:} \quad &\text{мин} = 0{,}07, \quad \text{второй} = 0{,}10, \quad \text{штраф} = 0{,}03,\\
\text{Столбец 2:} \quad &\text{мин} = 0{,}08, \quad \text{второй} = 0{,}09, \quad \text{штраф} = 0{,}01,\\
\text{Столбец 3:} \quad &\text{мин} = 0{,}06, \quad \text{второй} = 0{,}08, \quad \text{штраф} = 0{,}02,\\
\text{Столбец 4:} \quad &\text{мин} = 0{,}08, \quad \text{второй} = 0{,}09, \quad \text{штраф} = 0{,}01.
\end{align*}

Максимальный штраф: $0{,}03$ (строка 4, столбец 1). В строке 4 минимум $c_{42} = 0{,}08$.
Перевозим $x_{42} = \min(130, 145) = 130$. Остаток: $b_2 = 15$. Строка 4 исключена.

\textbf{Итерация 2.}

Штрафы по оставшимся строкам:
\begin{align*}
\text{Строка 1:} \quad &0{,}02,\\
\text{Строка 2:} \quad &0{,}02,\\
\text{Строка 3:} \quad &0{,}02.
\end{align*}

Штрафы по столбцам:
\begin{align*}
\text{Столбец 1:} \quad &\text{мин} = 0{,}07, \quad \text{второй} = 0{,}10, \quad \text{штраф} = 0{,}03,\\
\text{Столбец 2:} \quad &\text{мин} = 0{,}09, \quad \text{второй} = 0{,}14, \quad \text{штраф} = 0{,}05,\\
\text{Столбец 3:} \quad &\text{мин} = 0{,}06, \quad \text{второй} = 0{,}08, \quad \text{штраф} = 0{,}02,\\
\text{Столбец 4:} \quad &\text{мин} = 0{,}09, \quad \text{второй} = 0{,}10, \quad \text{штраф} = 0{,}01.
\end{align*}

Максимальный штраф: $0{,}05$ (столбец 2). В столбце 2 минимум $c_{12} = 0{,}09$.
Перевозим $x_{12} = \min(140, 15) = 15$. Остаток: $a_1 = 125$. Столбец 2 исключен.

\textbf{Итерация 3.}

Штрафы по строкам:
\begin{align*}
\text{Строка 1:} \quad &\text{мин} = 0{,}07, \quad \text{второй} = 0{,}12, \quad \text{штраф} = 0{,}05,\\
\text{Строка 2:} \quad &\text{мин} = 0{,}06, \quad \text{второй} = 0{,}08, \quad \text{штраф} = 0{,}02,\\
\text{Строка 3:} \quad &\text{мин} = 0{,}07, \quad \text{второй} = 0{,}10, \quad \text{штраф} = 0{,}03.
\end{align*}

Штрафы по столбцам:
\begin{align*}
\text{Столбец 1:} \quad &\text{мин} = 0{,}07, \quad \text{второй} = 0{,}10, \quad \text{штраф} = 0{,}03,\\
\text{Столбец 3:} \quad &\text{мин} = 0{,}06, \quad \text{второй} = 0{,}12, \quad \text{штраф} = 0{,}06,\\
\text{Столбец 4:} \quad &\text{мин} = 0{,}09, \quad \text{второй} = 0{,}10, \quad \text{штраф} = 0{,}01.
\end{align*}

Максимальный штраф: $0{,}06$ (столбец 3). В столбце 3 минимум $c_{23} = 0{,}06$.
Перевозим $x_{23} = \min(125, 155) = 125$. Остаток: $b_3 = 30$. Строка 2 исключена.

\textbf{Итерация 4.}

Штрафы по строкам:
\begin{align*}
\text{Строка 1:} \quad &\text{мин} = 0{,}07, \quad \text{второй} = 0{,}12, \quad \text{штраф} = 0{,}05,\\
\text{Строка 3:} \quad &\text{мин} = 0{,}07, \quad \text{второй} = 0{,}10, \quad \text{штраф} = 0{,}03.
\end{align*}

Штрафы по столбцам:
\begin{align*}
\text{Столбец 1:} \quad &\text{мин} = 0{,}07, \quad \text{второй} = 0{,}10, \quad \text{штраф} = 0{,}03,\\
\text{Столбец 3:} \quad &\text{мин} = 0{,}07, \quad \text{второй} = 0{,}12, \quad \text{штраф} = 0{,}05,\\
\text{Столбец 4:} \quad &\text{мин} = 0{,}09, \quad \text{второй} = 0{,}10, \quad \text{штраф} = 0{,}01.
\end{align*}

Максимальный штраф: $0{,}05$ (строка 1, столбец 3 одновременно). Выбираем $c_{13} = 0{,}12$.
Перевозим $x_{13} = \min(125, 30) = 30$. Остаток: $a_1 = 95$. Столбец 3 исключен.

\textbf{Итерация 5.}

Остались строки 1, 3 и столбцы 1, 4.
Минимум: $c_{11} = 0{,}07$. Перевозим $x_{11} = \min(95, 130) = 95$.
Остаток: $b_1 = 35$. Строка 1 исключена.

\textbf{Итерация 6.}

Остаток: $a_3 = 160$, нужно распределить между $b_1 = 35$ и $b_4 = 125$.
Перевозим $x_{31} = 35$ и $x_{34} = 125$.

\subsection{Матрица перевозок (метод Фогеля)}

\begin{table}[h!]
\centering
\caption{Распределение грузов методом Фогеля}
\label{tab:vogel}
\renewcommand{\arraystretch}{1.3}
\begin{tabular}{|c|c|c|c|c|c|}
\hline
 & $j=1$ & $j=2$ & $j=3$ & $j=4$ & $a_i$ \\
\hline
$i=1$ & $95$ & $15$ & $30$ & $0$ & $140$ \\
      & $(0{,}07)$ & $(0{,}09)$ & $(0{,}12)$ &  &  \\
\hline
$i=2$ & $0$ & $0$ & $125$ & $0$ & $125$ \\
      &  &  & $(0{,}06)$ &  &  \\
\hline
$i=3$ & $35$ & $0$ & $0$ & $125$ & $160$ \\
      & $(0{,}10)$ &  &  & $(0{,}09)$ &  \\
\hline
$i=4$ & $0$ & $130$ & $0$ & $0$ & $130$ \\
      &  & $(0{,}08)$ &  &  &  \\
\hline
$b_j$ & $130$ & $145$ & $155$ & $125$ & $555$ \\
\hline
\end{tabular}
\end{table}

\subsection{Расчёт стоимости}

Базисные переменные (ненулевые клетки):
\begin{align*}
x_{11} &= 95, \quad \text{стоимость} = 95 \cdot 0{,}07 = 6{,}65,\\
x_{12} &= 15, \quad \text{стоимость} = 15 \cdot 0{,}09 = 1{,}35,\\
x_{13} &= 30, \quad \text{стоимость} = 30 \cdot 0{,}12 = 3{,}60,\\
x_{23} &= 125, \quad \text{стоимость} = 125 \cdot 0{,}06 = 7{,}50,\\
x_{31} &= 35, \quad \text{стоимость} = 35 \cdot 0{,}10 = 3{,}50,\\
x_{34} &= 125, \quad \text{стоимость} = 125 \cdot 0{,}09 = 11{,}25,\\
x_{42} &= 130, \quad \text{стоимость} = 130 \cdot 0{,}08 = 10{,}40.
\end{align*}

\textbf{Общая стоимость метода Фогеля:}
\[
Z_{\text{Ф}} = 6{,}65 + 1{,}35 + 3{,}60 + 7{,}50 + 3{,}50 + 11{,}25 + 10{,}40 = \mathbf{44{,}25}.
\]

Число базисных переменных: $7 = m + n - 1 = 4 + 4 - 1 = 7$. ✓

\newpage

\section{Анализ результатов}

\subsection{Сравнение методов}

\begin{table}[h!]
\centering
\caption{Сравнение результатов методов}
\label{tab:methods_compare}
\renewcommand{\arraystretch}{1.4}
\begin{tabular}{|l|c|c|}
\hline
Метод & Общая стоимость & Базисные переменные \\
\hline
Минимального элемента & 44,25 & 7 \\
\hline
Фогеля & 44,25 & 7 \\
\hline
\end{tabular}
\end{table}

\subsection{Выводы}

\begin{enumerate}
\item Оба метода дали одинаковое распределение грузов и одинаковую стоимость $Z = 44{,}25$.

\item Это произошло потому, что матрица стоимостей имеет специальную структуру, в которой
оба алгоритма сходятся к одному решению.

\item \textbf{Метод минимального элемента:}
    \begin{itemize}
        \item Простой в реализации.
        \item Работает за $O(m \cdot n \cdot \log(m \cdot n))$ операций.
        \item Часто даёт неплохие начальные решения для небольших задач.
        \item Не гарантирует оптимальность.
    \end{itemize}

\item \textbf{Метод Фогеля:}
    \begin{itemize}
        \item Более сложный в реализации (требует расчёта штрафов на каждом шаге).
        \item Работает за $O(m^2 \cdot n^2)$ операций в худшем случае.
        \item Часто даёт лучшие начальные решения, чем метод минимального элемента.
        \item Не гарантирует оптимальность, но требует меньше итераций улучшения.
    \end{itemize}

\item Полученное решение необходимо проверить на оптимальность с помощью метода потенциалов
или решить задачу симплекс-методом.

\item \textbf{Матрица перевозок обоих методов}:
    \[
    X = \begin{pmatrix}
    95 & 15 & 30 & 0 \\
    0  & 0  & 125 & 0 \\
    35 & 0  & 0  & 125 \\
    0  & 130 & 0  & 0
    \end{pmatrix}.
    \]

\item \textbf{Общие затраты}: $Z = 44{,}25$ условных единиц (при единице груза за ячейку).

\end{enumerate}

\newpage

\section*{Приложение}
\addcontentsline{toc}{section}{Приложение}

\begin{lstlisting}[caption={Решение транспортной задачи методами минимального элемента и Фогеля}, label=lst:code]
import numpy as np

# Matritsa stoimosti
C = np.array([
    [0.07, 0.09, 0.12, 0.10],
    [0.11, 0.06, 0.08, 0.13],
    [0.10, 0.14, 0.07, 0.09],
    [0.12, 0.08, 0.11, 0.08]
])

# Predlozhenie i spros
a = np.array([140, 125, 160, 130])
b = np.array([130, 145, 155, 125])

m, n = C.shape
print(f"Sumbalnaya zadacha: sum(a)={sum(a)}, sum(b)={sum(b)}")

# ===== Metod minimalnogo elementa =====
def min_element_method(C, a, b):
    X = np.zeros((len(a), len(b)))
    a_left = a.copy()
    b_left = b.copy()
    
    while np.sum(a_left) > 0:
        # Najti minimum
        mask = np.ones_like(C, dtype=bool)
        for i in range(len(a)):
            if a_left[i] == 0:
                mask[i, :] = False
        for j in range(len(b)):
            if b_left[j] == 0:
                mask[:, j] = False
        
        if not np.any(mask):
            break
        
        C_masked = C.copy()
        C_masked[~mask] = np.inf
        min_idx = np.unravel_index(np.argmin(C_masked), C_masked.shape)
        i, j = min_idx
        
        vol = min(a_left[i], b_left[j])
        X[i, j] = vol
        a_left[i] -= vol
        b_left[j] -= vol
    
    cost = np.sum(X * C)
    return X, cost

# ===== Metod Vogelya =====
def vogel_method(C, a, b):
    X = np.zeros((len(a), len(b)))
    a_left = a.copy()
    b_left = b.copy()
    
    while np.sum(a_left) > 0:
        active_rows = [i for i in range(len(a)) if a_left[i] > 0]
        active_cols = [j for j in range(len(b)) if b_left[j] > 0]
        
        if len(active_rows) == 0 or len(active_cols) == 0:
            break
        
        # Calculate penalties
        penalties_rows = []
        for i in active_rows:
            row = C[i, active_cols]
            if len(row) >= 2:
                row_sorted = np.sort(row)
                penalty = row_sorted[1] - row_sorted[0]
            else:
                penalty = 0
            penalties_rows.append((i, penalty))
        
        penalties_cols = []
        for j in active_cols:
            col = C[active_rows, j]
            if len(col) >= 2:
                col_sorted = np.sort(col)
                penalty = col_sorted[1] - col_sorted[0]
            else:
                penalty = 0
            penalties_cols.append((j, penalty))
        
        # Max penalty
        max_row_penalty = max(penalties_rows, key=lambda x: x[1])[1]
        max_col_penalty = max(penalties_cols, key=lambda x: x[1])[1]
        
        if max_row_penalty >= max_col_penalty:
            # Find min in row
            i = max(penalties_rows, key=lambda x: x[1])[0]
            j = active_cols[np.argmin(C[i, active_cols])]
        else:
            # Find min in col
            j = max(penalties_cols, key=lambda x: x[1])[0]
            i = active_rows[np.argmin(C[active_rows, j])]
        
        vol = min(a_left[i], b_left[j])
        X[i, j] = vol
        a_left[i] -= vol
        b_left[j] -= vol
    
    cost = np.sum(X * C)
    return X, cost

# Run both methods
X_min, cost_min = min_element_method(C, a, b)
X_vogel, cost_vogel = vogel_method(C, a, b)

print("\n===== Metod minimalnogo elementa =====")
print("Matritsa perevozok:\n", X_min)
print(f"Obshaya stoimost: {cost_min:.2f}")

print("\n===== Metod Vogelya =====")
print("Matritsa perevozok:\n", X_vogel)
print(f"Obshaya stoimost: {cost_vogel:.2f}")

print("\n===== Sravnenie =====")
print(f"Raznitsa stoimosti: {abs(cost_min - cost_vogel):.4f}")
\end{lstlisting}

\end{document}